\subsection{Assumptions, efficiencies, and formalization boundaries}\label{sec:dpr:interfaces}

\begin{definition}[The three proof-system layers]\label{def:dpr:layers}
The development distinguishes:
\begin{enumerate}[label=\textup{(\roman*)}]
\item \emph{Algebraic reduction:}
  exact packing, the tensor evaluation
  $S$, and the sumcheck polynomial $H$;
\item \emph{Committed evaluation:}
  a fixed packed polynomial $T'$,
  a binding representation $C'$, and
  an underlying large-field evaluation proof;
\item \emph{Zero-knowledge implementation:}
  commitments, masked messages and
  a simulator producing the full adaptive
  verifier view.
\end{enumerate}
Only the first layer is proved
unconditionally in this chapter.
The second has an explicit
conditional soundness theorem;
the third remains an obligation.
\end{definition}

\begin{lemma}[No increase in the packed Boolean table size]\label{lem:dpr:table-size}
The original Boolean table has
$2^{\ell+h}=n\,2^h$ binary entries.
The packed table has $2^h$
elements of $L$, each represented
by $n$ binary coordinates.
Thus both data representations
contain exactly $n\,2^h$
binary coordinates.
\end{lemma}
\begin{proof}
The set $I_\ell$ has cardinality
$2^\ell=n$ and $I_h$ has
cardinality $2^h$, so the original
table has $n2^h$ entries.
The packed table has one element
of $L$ for each $j\in I_h$.
Since $[L:\F_2]=n$, every field
element has exactly $n$ binary
basis coordinates, and hence
the packed table has $n2^h$
such coordinates.
\end{proof}

\begin{lemma}[Message types and sumcheck round count]\label{lem:dpr:round-count}
The tensor reduction sends one
initial $A$-element $S_0$,
then exactly $h$ univariate
$A$-polynomials of degree at most
$2$, and one final field
element $u\in L$ before
invoking the delegated
large-field evaluation protocol.
\end{lemma}
\begin{proof}
The number and type of messages
are specified in
Definition~\ref{def:dpr:full-protocol}:
one initial tensor message,
one polynomial in each iteration
$i=0,\ldots,h-1$, and one claimed
evaluation in the last reduction
step. The degree bound follows
from Definition~\ref{def:dpr:sumcheck}.
\end{proof}

\begin{remark}[Why this is not the masked protocol]
The initial tensor message
$S_0$ reveals the partial table
$s^*$ by Lemma~\ref{lem:dpr:leaks-partials}.
The masking proposal in
Chapters~6--7 changes the
public messages and proof
constraints; it cannot inherit
the original sumcheck theorem
merely by reusing the name
ring switching.
Its own proof needs an
end-to-end simulation of
$(m,y',U,A)$ and all
commitment/opening messages,
with every challenge timing
and correlation checked.
\end{remark}

\begin{definition}[Chapter 8 Lean certification contract]\label{def:dpr:lean-contract}
To certify the exact compiler
in Lean~4, the formal model
will provide:
\begin{enumerate}[label=\textup{(\roman*)}]
\item the finite extension
$L/\F_2$, two $L$-module
structures on $A$, tensor
product and its universal map;
\item binary packing/unpacking
and the polynomial $T'$;
\item the typed column functional
and tensor equality interpolation;
\item the polynomial ring $A[X]$
and its root bound restricted
to embedded units;
\item causal sumcheck messages,
independent challenge sampling,
and its soundness event; and
\item a separate abstract
large-field commitment interface
with an explicitly parameterized
evaluation error.
\end{enumerate}
No Lean declaration is claimed
certified until its code has
been compiled and checked.
\end{definition}

\begin{remark}[History and non-novelty]
The \emph{unmasked} compiler and
its tensor-sumcheck soundness
come from Diamond--Posen
\cite[Construction~4.1, Theorems~4.2 and~4.4]{DiamondPosen2024}.
The proof of the embedded-field
root bound corresponds to
their Lemma~4.5.
The present chapter reorganizes
those established arguments
as small dependency-ordered
proofs, specialized to the
Boolean-kernel basis already
used in this manuscript.
It is not presented as a
new ring-switching construction.
\end{remark}
